\documentclass[11pt]{article}
\usepackage{booktabs,xcolor,colortbl,caption,array,graphicx,amsmath}
\usepackage[landscape,margin=1cm]{geometry}
\begin{document}\small\setcounter{table}{7}
\begin{table}[!h]
\centering
\caption{\label{tab:tab:dm_decomp}Dark matter par classe d'actif sous deux scénarios, moy. 2013--2022 (\% du PIB)}
\centering
\fontsize{9.5}{11.5}\selectfont
\begin{tabular}[t]{lrrrrrrrrrr}
\toprule
\multicolumn{1}{c}{\textbf{ }} & \multicolumn{5}{c}{\textbf{Scénario A: taux G\&R universels}} & \multicolumn{5}{c}{\textbf{Scénario B: taux différenciés}} \\
\cmidrule(l{3pt}r{3pt}){2-6} \cmidrule(l{3pt}r{3pt}){7-11}
Entité & Total & FDI & Actions & Dette & Autres & Total & FDI & Actions & Dette & Autres\\
\midrule
China & NaN & NaN & NaN & NaN & NaN & NaN & NaN & NaN & NaN & NaN\\
European Union & -0.3 & -4.8 & 4.2 & 5.2 & -5.0 & -0.3 & -4.8 & 4.2 & 5.2 & -5.0\\
Japan & 10.5 & -15.2 & 6.7 & 41.4 & -22.3 & 10.5 & -15.2 & 6.7 & 41.4 & -22.3\\
United States & 44.0 & 18.3 & 7.8 & 19.0 & -1.1 & 26.4 & 16.7 & 6.8 & 4.0 & -1.2\\
\bottomrule
\multicolumn{11}{l}{\rule{0pt}{1em}\textit{\textit{Notes:} }}\\
\multicolumn{11}{l}{\rule{0pt}{1em}DM par composante = NFA implicite (NII/r) $-$ stock officiel (EWN). Une part FDI dominante valide le canal \textit{knowledge} de H\&S; une part dette dominante valide le canal actif s\^{u}r (Caballero, Farhi \& Gourinchas 2017). Pour les \'{e}conomies avancées non-privilège et l'UE, Scén. A = Scén. B.}\\
\end{tabular}
\end{table}
\end{document}
